% tex.web S367 makes \noexpand's token a `dont_expand' token (S358): `\relax'
% to whatever reads it next -- including the expander, so \csname is not run --
% and the same control sequence again once anything stores it. Reaching
% \noexpand through \expandafter is the case that hands the marked token to a
% scanner other than the one that saw \noexpand itself.
\catcode`\{=1 \catcode`\}=2 \catcode`\#=6
\errorcontextlines=5
\def\a{A}
\message{[\expandafter\noexpand\csname relax\endcsname]}
\message{[\expandafter\noexpand\expandafter\a\noexpand\a\noexpand\csname x\endcsname]}
\message{[\ifx\noexpand\a\noexpand\csname\fi][\ifx\noexpand\a\relax T\else F\fi]}
\message{[\if\noexpand\a\relax T\else F\fi\ifcat\noexpand\a\relax T\else F\fi]}
\expandafter\def\expandafter\b\expandafter{\noexpand\a}
\message{[\meaning\b]}
\edef\c{\expandafter\noexpand\a\noexpand\a}
\message{[\meaning\c]}
\expandafter\show\noexpand\a
\message{[end]\noexpand}
\end
